\documentclass{article}
\usepackage[T1]{fontenc}
\usepackage[french]{babel}
\usepackage[french, ruled, vlined]{algorithm2e}
\usepackage{amsmath, amssymb}
\usepackage{float}

\title{TER}
\author{Alexandre François}
\date{February 2026}

\begin{document}

    \maketitle

    \section{Recherche Exhaustive par Différence d'Entrée (Algorithme Naïf Optimisé)}\label{sec:recherche-exhaustive-par-difference-d'entree-(algorithme-naif-optimise)}

    \subsection{Principe et Algorithme}\label{subsec:principe-et-algorithme}

    \noindent \textbf{Problèmatique.} L'approche naïve, qui consisterait à tester toutes les paires de textes clairs $(x,y)$ possibles, possède une complexité prohibitive de $O(2^n)$. Pour un bloc de taille $n = 16$, cela représente déjà 4 milliards de tests. Pour $n = 64$, c'est physiquement impossible.

    \vspace{0.2cm}

    \noindent \textbf{Stratégie.} L'algorithme contourne l'explosion combinatoire en structurant la recherche autour des différences d'entrée $\alpha$. Au lieu de tester l'ensemble de l'espace, nous procédons par échantillonnage : pour chaque $\alpha$, nous générons un volume restreint de $c/p$ paires de messages. Cette approche permet de vérifier si une différence d'entrée "active" une différence de sortie $\beta$ particulière. Si une telle paire $(x, x \oplus \alpha)$ émerge du bruit de fond, elle confirme l'existence d'une faille sans qu'il soit nécessaire de parcourir les milliards de combinaisons possibles.

    \vspace{0.2cm}

    \noindent \textbf{Complexité.} Cette approche présente une complexité globale de $O(2^n \cdot \frac{1}{p})$. Elle reste ainsi envisageable pour des tailles de blocs modestes (par exemple jusqu'à $n = 32$ avec une puissance de calcul suffisante), mais devient impraticable pour des tailles standards comme $n = 64$.

    \begin{algorithm}[H]
        \caption{Recherche Différentielle Standard (Naïve Optimisée)}
        \SetKwInOut{Input}{Entrée}
        \SetKwInOut{Output}{Sortie}

        \Input{Taille de bloc $n$, Probabilité cible $p$, Constante de succès $c$.}
        \Output{Affichage des différentielles candidates $(\alpha, \beta)$ significatives.}

        \For{toute différence d'entrée non nulle $\alpha \in \{0,1\}^n \setminus \{0\}$}{
            Initialiser les compteurs $Count_\beta \leftarrow 0$\;

            \tcp{Phase d'échantillonnage}
            \For{$k=1$ \KwTo $\lceil c / p \rceil$}{
                Choisir aléatoirement $x \in \{0,1\}^n$\;
                Calculer $\beta = E(x) \oplus E(x \oplus \alpha)$\;
                $Count_\beta \leftarrow Count_\beta + 1$\;
            }

            \tcp{Phase d'analyse}
            \ForEach{$\beta$ observé tel que $Count_\beta > 1$}{
                Afficher $(\alpha, \beta)$;
            }
        }
    \end{algorithm}

    \subsection{Analyse Statistique et Dimensionnement}

    La fiabilité de cet algorithme ne repose pas sur le hasard, mais sur un dimensionnement rigoureux de l'échantillon $N$. Nous devons distinguer une véritable faille cryptographique (le Signal) des coïncidences fortuites (le Bruit). Cette section démontre pourquoi un choix naïf de $N$ conduit à l'échec et justifie l'introduction d'un facteur de sécurité $c$.

    \subsubsection{Les limites de l'approche unitaire ($N=1/p$)}

    Intuitivement, pour détecter une caractéristique de probabilité $p$, on pourrait penser qu'il suffit de tester $N = 1/p$ paires, espérant ainsi obtenir en moyenne une occurrence. Cependant, cette intuition est trompeuse en cryptanalyse.

    Modélisons le nombre d'apparitions de la différentielle par une variable aléatoire $X$. Dans le cas d'un chiffrement aléatoire, ce processus de comptage suit une loi binomiale $B(N, P_{rand})$.
    Puisque la probabilité $P_{rand} = 2^{-n}$ est infime et que le nombre d'essais $N$ est grand, nous pouvons appliquer le Théorème des événements rares pour approximer cette distribution par une loi de Poisson de paramètre $\lambda = N \cdot 2^{-n}$.
    Cette modélisation est fondamentale pour établir le rapport Signal/Bruit (S/N) qui détermine la complexité de l'attaque\footnote{Le concept de Rapport Signal/Bruit en cryptanalyse différentielle a été introduit par E. Biham et A. Shamir dans \textit{"Differential Cryptanalysis of DES-like Cryptosystems"} (1991). Ils y établissent que le compteur de la bonne clé doit se distinguer statistiquement du bruit de fond généré par les mauvaises clés.}.

    Si nous fixons $N=1/p$, l'espérance est effectivement de $\lambda = 1$. Or, la probabilité mathématique de ne jamais observer la différentielle (cas $k=0$) est loin d'être négligeable :
    $$P(X = 0) = e^{-1} \approx 36.8\%$$
    Cela signifie qu'une fois sur trois, l'attaque échoue simplement par manque de chance.

    Le problème s'aggrave lorsque l'on considère le filtrage du bruit. Pour éviter les faux positifs, notre algorithme ignore les événements uniques ($Count=1$). L'attaque échoue donc si la différentielle apparaît 0 fois ou 1 seule fois. La probabilité d'échec cumulée devient critique :
    $$P(\text{Échec}) = P(X=0) + P(X=1) = e^{-1} + e^{-1} \approx 73.6\%$$
    Avec un taux d'échec supérieur à 70\%, l'approche unitaire est inutilisable.



    \subsubsection{Optimisation via le facteur de sécurité ($N=c/p$)}

    Pour pallier ce défaut, nous introduisons un facteur de sécurité $c$ (avec $c > 1$), en fixant la taille de l'échantillon à $N = c/p$. L'espérance du nombre de collisions devient alors $\lambda = c$.

    Nous cherchons la valeur de $c$ qui maximise la probabilité que la différentielle apparaisse au moins 2 fois (condition de succès). Cette probabilité est donnée par :
    $$ P(\text{Succès}) = P(X \ge 2) = 1 - e^{-c}(1 + c) $$

    L'analyse de cette fonction montre une progression rapide de la fiabilité :

    \begin{table}[H]
        \centering
        \begin{tabular}{|c|c|c|c|}
            \hline
            \textbf{Facteur ($c$)} & \textbf{Espérance ($\lambda$)} & \textbf{Probabilité d'Échec} & \textbf{Probabilité de Succès} \\
            \hline
            $c=1$ & 1 & 73.6\% & 26.4\% \\
            \hline
            $c=2$ & 2 & 40.6\% & 59.4\% \\
            \hline
            $c=3$ & 3 & 19.9\% & \textbf{80.1\%} \\
            \hline
            $c=4$ & 4 & 9.2\% & \textbf{90.8\%} \\
            \hline
            $c=5$ & 5 & 4.0\% & \textbf{96.0\%} \\
            \hline
        \end{tabular}
        \caption{Impact du facteur $c$ sur la détection de la faille.}
        \label{tab:poisson_c}
    \end{table}

    Il apparaît clairement qu'un facteur $c \ge 4$ est nécessaire pour garantir un taux de détection supérieur à 90\%, rendant l'attaque statistiquement robuste.

    \subsubsection{Validation finale : Rapport Signal/Bruit}

    Augmenter $N$ améliore la détection du signal, mais augmente aussi mécaniquement le nombre de fausses collisions dues au hasard. Il faut donc vérifier que notre seuil de filtrage ($Count > 1$) reste capable de discriminer le vrai du faux.

    Prenons l'exemple d'un chiffrement SPN ($n=16$) avec une faille cible $p=2^{-10}$ et un facteur $c=4$. Nous comparons deux hypothèses :
    \begin{itemize}
        \item \textbf{Hypothèse du Signal (La faille) :} L'espérance est $\lambda_S = c = 4$.
        \item \textbf{Hypothèse du Bruit (Le hasard) :} L'espérance est $\lambda_B = N \cdot 2^{-n} = \frac{4}{2^{-10} \cdot 2^{16}} = 2^{-4} = 0.0625$.
    \end{itemize}

    La probabilité de franchir le seuil de 2 collisions est radicalement différente pour les deux cas :
    $$ P(Signal \ge 2) \approx \mathbf{90.8\%} $$
    $$ P(Bruit \ge 2) \approx \mathbf{0.19\%} $$

    \textbf{Conclusion :} Le critère choisi agit comme un filtre passe-haut très sélectif. Il capture la quasi-totalité des différentielles réelles tout en rejetant 99.8\% du bruit aléatoire.

\end{document}
